\documentclass[../main.tex]{subfiles}
\graphicspath{{\subfix{../images/}}}
\begin{document}






\begin{example}
    Nekonečně rozměrný dynamický systém:
    \begin{equation}
        \partial_t u = D\partial^2_{xx} u, v (0,\infty) \times (a,b)
    \end{equation}

    Okrajové podmínky $u|_{\bar{x}=a} = 0, u|_{\bar{x}=b} = 0$ a počáteční podmínky $u|_{t=0} = u_{ini}(\bar{x})$ v $(a,b)$


    ÚKOL: zapsat analytické řešení ve tvaru řady a určit, jaký DS generuje.   

    Pro usnadnění provedeme transformaci $x=\bar{x}-a$
    Transformované okrajové podmínky jsou $u|_{x=0} = 0, u|_{x=b-a} = 0$

    Řešíme pomoci separace proměnných. 

    \begin{equation}
        u(x,t) = S(t)T(x)
    \end{equation}
    Rovnice je pak 

    \begin{equation}
        T(x) \frac{d S(t)}{dt} = \frac{d^2 T(x)}{dx^2} S(t)
    \end{equation}


    Upravíme
    \begin{equation}
        \frac{T(x)}{\frac{d^2 T(x)}{dx^2}} = \frac{S(t)}{\frac{d S(t)}{dt}}
    \end{equation}

    Levá strana závisí pouze na x, levá pouze na t, proto se obě strany rovnají nějaké konstantě C.

    Z pravé strany máme 
    \begin{equation}
        \frac{d S(t)}{dt} = CS(t)
    \end{equation}

    Z toho dostáváme
    \begin{equation}
        S(t) = A e^{C t}
    \end{equation}

    Z levé strany máme
    \begin{equation}
        \frac{d^2 T(x)}{dx^2} = C T(x)
    \end{equation}

    Řešením je
    \begin{equation}
        T(x) = b_1 e^{\sqrt{C}x} + b_2 e^{-\sqrt{C}x}
    \end{equation}

    Z okrajových podmínek víme, že $T(a) = T(b) = 0$.

    Pokud je $C\geq 0$, úloha má řešení pouze pro $u=0$. Proto $C < 0$. Přeznačme $\alpha := -C$

    Z toho dostáváme 
    \begin{equation}
        T(x) = a_1 \sin(\sqrt{\alpha}x) + a_2 \cos(\sqrt{\alpha}x)
    \end{equation}

    Dosadíme okrajovou podmínku $u|_{x=0} = 0$ a máme 
    \begin{equation}
        T(x) = a_1 \sin(\sqrt{\alpha}x)
    \end{equation}

    Z druhé okrajové podmínky máme
    \begin{equation}
        T(x) = a_1 \sin(n\frac{\pi}{b-a}x)
    \end{equation},
    kde $n$ je celé číslo.
 
    Tedy $\alpha = n\frac{\pi}{b-a}$

    Celkově dostáváme řešení módů: 

    \begin{equation}
        u(x,t) = \sum_{n=1}^{\infty} C_n \sin(n\frac{\pi}{b-a}x) \exp(-\frac{n^2\pi^2t}{(b-a)^2})
    \end{equation}
    kde koeficienty $C_n$ se dostaly rozvojem počáteční podmínky do Fourierovy řady.

    Proveďme zpětnou transformaci $\bar{x}=x+a$
    
    \begin{equation}
        u(\bar{x},t) = \sum_{n=1}^{\infty} C_n \sin(n\frac{\pi}{b-a} (\bar{x} - a)) \exp(-\frac{n^2\pi^2t}{(b-a)^2})
    \end{equation}

    Řada diverguje pro $t<0$, dynamický systém je neinvertovatelný.
\end{example}











































\end{document}